\documentclass[12pt,a4paper]{article}
\usepackage[utf8]{inputenc}
\usepackage[T1]{fontenc}
\usepackage[italian]{babel}
\usepackage{libertinus}
\usepackage{amssymb}
\usepackage[a4paper,top=25mm,bottom=25mm,left=22mm,right=22mm,headheight=15pt]{geometry}
\usepackage{graphicx}
\graphicspath{{../../../assets/}{../../assets/}{../assets/}{assets/}}
\usepackage[table]{xcolor}
\usepackage{tabularx,array}
\usepackage{titlesec}
\usepackage{titletoc}
\usepackage{microtype}
\usepackage{fancyhdr}
\usepackage{enumitem}
\usepackage[hidelinks]{hyperref}

% ---------- Dati del verbale (da impostare una sola volta) ----------
\newcommand{\vnum}{9}                     % numero progressivo del verbale interno
\newcommand{\vdata}{2026-11-19}           % data della riunione, formato ISO
\newcommand{\vdataestesa}{19 novembre 2026}   % data della riunione per esteso (es. 19 novembre 2026)

% Palette colori del gruppo Bending Forks
\definecolor{bfblue}{HTML}{1B3F9E}
\definecolor{bfrow}{HTML}{F2F5FC}
\definecolor{bfline}{HTML}{C9D1E6}
\definecolor{bfhead}{HTML}{2E6BE6}

% Macro di formattazione interna Bending Forks
\newcommand{\ph}[1]{\textcolor{red}{[#1]}}                        % segnaposto da completare
\newcommand{\gl}[1]{\emph{#1}\textsubscript{\scriptsize G}}       % termine del Glossario
\newcommand{\thc}[1]{\textcolor{white}{\textbf{#1}}}              % cella di intestazione delle tabelle
\newcommand{\rl}{\arrayrulecolor{bfline}\hline}                   % riga di separazione delle tabelle
\newcommand{\presente}{$\blacksquare$ & $\square$}                % colonne Presente / Assente
\newcommand{\assente}{$\square$ & $\blacksquare$}
\newcommand{\tightlist}{\setlength{\itemsep}{1pt}\setlength{\parskip}{0pt}}

\setlist{topsep=2pt}
\setlist[itemize]{leftmargin=1.3em,label=\textbullet}
\setlist[enumerate]{leftmargin=1.6em}
\setlength{\parindent}{0pt}
\sloppy
\setlength{\emergencystretch}{2em}
\setlength{\parskip}{5pt}
\renewcommand{\arraystretch}{1.45}
\setcounter{secnumdepth}{3}
\setcounter{tocdepth}{3}

% Formattazione titoli sezioni
\titleformat{\section}{\LARGE\bfseries\color{bfblue}\raggedright}{\thesection.}{0.3em}{}[{\color{bfblue}\vspace{-2pt}\titlerule[0.9pt]}]
\titleformat{\subsection}{\Large\bfseries\color{bfblue}\raggedright}{\thesubsection.}{0.3em}{}
\titleformat{\subsubsection}{\normalsize\bfseries\raggedright}{\thesubsubsection.}{0.3em}{}
\titlespacing*{\section}{0pt}{14pt}{6pt}
\titlespacing*{\subsection}{0pt}{10pt}{4pt}
\titlespacing*{\subsubsection}{0pt}{10pt}{3pt}

% Indice cliccabile con trattini fino al numero di pagina
\contentsmargin{2.2em}
\newcommand{\bfleader}{\titlerule*[0.32em]{\textendash}}
\titlecontents{section}[0pt]{\addvspace{5pt}\large\bfseries}{\thecontentslabel.\ }{}{\hspace{0.4em}\bfleader\contentspage}
\titlecontents{subsection}[1.2em]{\addvspace{0pt}}{\thecontentslabel.\ }{}{\hspace{0.4em}\bfleader\contentspage}
\titlecontents{subsubsection}[2.4em]{\addvspace{0pt}}{\thecontentslabel.\ }{}{\hspace{0.4em}\bfleader\contentspage}

% Intestazione e piè di pagina
\pagestyle{fancy}
\fancyhf{}
\fancyhead[L]{\footnotesize\sffamily\bfseries\color{bfhead}\textls[120]{BENDING FORKS}}
\fancyhead[R]{\footnotesize\sffamily Verbale interno n. \vnum}
\fancyfoot[C]{\small Pagina \thepage}
\renewcommand{\headrulewidth}{0.9pt}

% Colorazione alternata delle righe delle tabelle
\newcommand{\bftable}{\rowcolors{2}{bfrow}{white}}
\usepackage{etoolbox}
\AtBeginEnvironment{tabularx}{\bftable}

\begin{document}

% ---------- Copertina ----------
\begin{titlepage}
\centering
\vspace*{4mm}
\includegraphics[width=105mm]{logo-bending-forks-trasparente.png}\\[4mm]
\underline{bending.forks.swe@gmail.com}

\vspace{42mm}
{\fontsize{24}{28}\selectfont\bfseries Verbale interno n. \vnum\par}
\vspace{4mm}
{\Large\color{bfblue} Riunione del \vdataestesa\par}

\vfill
\begin{tabular}[t]{r|l}
\textbf{Responsabile} & Filippo Rampazzo \\
\textbf{Segretario} & Lorenzo Milanese \\
\textbf{Verificatore} & Nicholas Salmaso \\
\end{tabular}
\hspace{14mm}
\begin{tabular}[t]{r|l}
\textbf{Versione} & 0.1.1 \\
\textbf{Tipo} & Verbale interno \\
\textbf{Destinatari} & \emph{Gruppo Bending Forks} \\
 & \emph{Prof. Tullio Vardanega} \\
 & \emph{Prof. Riccardo Cardin} \\
\end{tabular}

\vspace{12mm}
{\large\emph{\textbf{Abstract dei contenuti:}}}\\[1mm]
\begin{minipage}{150mm}\centering Sintesi e decisioni emerse durante la riunione interna del gruppo del \vdataestesa, riguardante la revisione della bozza dell'Analisi dei requisiti e dei casi d'uso del progetto, le scelte sulla comunicazione tra dispositivo e server, la preparazione della riunione con la proponente del 20 novembre 2026 e alcune regole sull'uso di GitHub.\end{minipage}
\vspace*{6mm}
\end{titlepage}

% ---------- Indice ----------
\vspace*{2mm}
{\fontsize{22}{26}\selectfont\bfseries Indice\par}
\vspace{6mm}
\makeatletter\@starttoc{toc}\makeatother
\newpage

% ---------- Registro delle versioni ----------
% Una riga per versione, dalla più recente alla più vecchia.
% Il Verificatore inserisce il proprio nome nella riga della versione che ha verificato.
\hypertarget{registro-delle-versioni}{%
\section{Registro delle versioni}}

\begin{center}
\begin{tabularx}{\linewidth}{
  >{\hsize=0.55\hsize\raggedright\arraybackslash}X
  >{\hsize=0.75\hsize\raggedright\arraybackslash}X
  >{\hsize=0.85\hsize\raggedright\arraybackslash}X
  >{\hsize=0.85\hsize\raggedright\arraybackslash}X
  >{\hsize=2.00\hsize\raggedright\arraybackslash}X
}
\rowcolor{bfblue}\thc{Versione} & \thc{Data} & \thc{Autore} & \thc{Verificatore} & \thc{Descrizione} \\
0.1.1 & 2026-11-23 & Lorenzo Milanese & Nicholas Salmaso & Correzioni della verifica: segnati i termini del Glossario e aggiunte le firme \\ \rl
0.1.0 & 2026-11-20 & Lorenzo Milanese & Nicholas Salmaso & Redazione del verbale \\ \rl
\end{tabularx}
\end{center}

% ---------- Informazioni generali ----------
\hypertarget{informazioni-generali}{%
\section{Informazioni generali}}

\hypertarget{durata-e-luogo}{%
\subsection{Durata e luogo}}

\begin{itemize}
\tightlist
\item \textbf{Data:} \vdata;
\item \textbf{Inizio:} 15:30;
\item \textbf{Fine:} 18:15;
\item \textbf{Luogo / modalità:} in remoto, tramite Discord.
\end{itemize}

\hypertarget{partecipanti}{%
\subsection{Partecipanti}}

% Per ogni componente: ruolo in carica nel periodo e \presente oppure \assente.
\begin{center}
\begin{tabularx}{\linewidth}{>{\raggedright\arraybackslash}X>{\centering\arraybackslash}X>{\centering\arraybackslash}p{2.2cm}>{\centering\arraybackslash}p{2.2cm}}
\rowcolor{bfblue}\thc{Nome e cognome} & \thc{Ruolo} & \thc{Presente} & \thc{Assente} \\
Marco Beltrame & Progettista & \presente \\ \rl
Lorenzo Milanese & Amministratore & \presente \\ \rl
Filippo Rampazzo & Responsabile & \presente \\ \rl
Nicholas Salmaso & Verificatore & \presente \\ \rl
Luca Stangherlin & Analista & \presente \\ \rl
Aldo Toffoli & Progettista & \presente \\ \rl
Matteo Wu & Analista & \presente \\ \rl
\end{tabularx}
\end{center}

% ---------- Ordine del giorno ----------
\hypertarget{ordine-del-giorno}{%
\section{Ordine del giorno}}

\begin{enumerate}
\tightlist
\item \hyperlink{odg-1}{Revisione dei requisiti funzionali dell'Analisi dei requisiti};
\item \hyperlink{odg-2}{Comunicazione tra dispositivo e server};
\item \hyperlink{odg-3}{Associazione e configurazione iniziale del dispositivo};
\item \hyperlink{odg-4}{Requisiti di qualità, prestazionali e di vincolo};
\item \hyperlink{odg-5}{Diagrammi, casi d'uso e unione delle due bozze};
\item \hyperlink{odg-6}{Preparazione della riunione con la proponente del 20 novembre 2026};
\item \hyperlink{varie-ed-eventuali}{Varie ed eventuali}.
\end{enumerate}
\newpage

% ---------- Discussione e decisioni ----------
% Una sottosezione per ogni punto dell'ordine del giorno, con lo stesso titolo usato nell'elenco.
% Ogni decisione ha il codice [VI-<n>-<nn>]: <n> è il numero del verbale (\vnum),
% <nn> il numero progressivo della decisione nel verbale, a due cifre (01, 02, ...).
% Motivazione, incaricati e scadenza si scrivono solo se il punto contiene decisioni;
% un punto senza decisioni ha solo la discussione.
\hypertarget{discussione-e-decisioni}{%
\section{Discussione e decisioni}}

\hypertarget{odg-1}{%
\subsection{Revisione dei requisiti funzionali dell'Analisi dei requisiti}}
\textbf{Discussione:}\\
Il gruppo rilegge la bozza dell'\gl{Analisi dei requisiti} di Beltrame e Toffoli, giudicata una buona base; vanno controllati alcuni salti nella numerazione dei requisiti. Principali osservazioni:
\begin{itemize}
\tightlist
\item \textbf{Chiamate tra dispositivi}: i requisiti da RF62 a RF66 sono ritenuti poco utili; la rimozione va confermata con la proponente;
\item \textbf{Account}: l'autenticazione è necessaria anche con server locale, perché la web app è esposta in rete;
\item \textbf{\emph{Provider} AI}: il supporto a un modello locale va chiarito con la proponente;
\item \textbf{Lingua}: va verificato il supporto multilingua di Whisper e Piper;
\item \textbf{Server \gl{MCP}}: la dashboard della web app verifica lo stato dei server solo quando viene aperta.
\end{itemize}

\textbf{Decisioni prese:}
\begin{itemize}
\tightlist
\item[$\checkmark$] \textbf{[VI-\vnum-01]} Rimozione del requisito RF142 sull'integrazione con dispositivi di altri ecosistemi;
\item[$\checkmark$] \textbf{[VI-\vnum-02]} Requisito RF61 classificato come desiderabile;
\item[$\checkmark$] \textbf{[VI-\vnum-03]} Ogni utente, al momento della registrazione, riceve automaticamente una casa. Può condividerla con altri utenti assegnando a ciascuno il ruolo di amministratore oppure di membro, che può solo visualizzare. Non si introducono famiglie o gruppi;
\item[$\checkmark$] \textbf{[VI-\vnum-04]} Il \emph{provider} AI è configurabile per singolo dispositivo; per l'\gl{MVP} è sufficiente un \emph{provider} remoto tramite chiave \gl{API};
\item[$\checkmark$] \textbf{[VI-\vnum-05]} Il dispositivo supporterà con certezza la lingua inglese; il supporto all'italiano è ancora da verificare, perché dipende dagli strumenti di riconoscimento e sintesi vocale che verranno scelti;
\item[$\checkmark$] \textbf{[VI-\vnum-06]} \gl{Wake word} personalizzata classificata come desiderabile;
\item[$\checkmark$] \textbf{[VI-\vnum-07]} Riscontro vocale dopo ogni azione e conferma a voce prima delle azioni con effetti esterni sono obbligatori;
\item[$\checkmark$] \textbf{[VI-\vnum-08]} Gli errori dei server MCP sono segnalati al momento della richiesta, senza \emph{polling};
\item[$\checkmark$] \textbf{[VI-\vnum-09]} RF167: visualizzazione e modifica del contesto del dispositivo riservate all'amministratore.
\end{itemize}

\textbf{Incaricati:} Beltrame, Toffoli, Stangherlin, Wu\\
\textbf{Scadenza:} 2026-11-27

\hypertarget{odg-2}{%
\subsection{Comunicazione tra dispositivo e server}}
\textbf{Discussione:}\\
Il server non può contattare il dispositivo di propria iniziativa: le alternative sono un \emph{polling} periodico o una connessione WebSocket sempre aperta, e si sceglie la seconda. L'audio viene inviato solo dopo la \emph{wake word}, riconosciuta in locale, e la risposta del modello arriva in \emph{streaming}.

\textbf{Decisioni prese:}
\begin{itemize}
\tightlist
\item[$\checkmark$] \textbf{[VI-\vnum-10]} La connessione WebSocket tra dispositivo e server resta sempre aperta e funge da canale di controllo;
\item[$\checkmark$] \textbf{[VI-\vnum-11]} Il server notifica gli aggiornamenti del \emph{firmware} tramite la connessione WebSocket; il dispositivo comunica la propria versione all'avvio e a ogni riconnessione;
\item[$\checkmark$] \textbf{[VI-\vnum-12]} Il \emph{token} del dispositivo è salvato nella memoria \emph{flash} e ha una scadenza, da definire in progettazione.
\end{itemize}

\textbf{Motivazione:}\\
La connessione sempre aperta rende immediati aggiornamenti e notifiche dal server e riduce la latenza della prima risposta, con un consumo di risorse trascurabile.

\hypertarget{odg-3}{%
\subsection{Associazione e configurazione iniziale del dispositivo}}
\textbf{Discussione:}\\
Non si raggiunge una conclusione sulla necessità del codice monouso per l'associazione (RF43, RF44), dato che il dispositivo non ha schermo né tastiera. Per la prima configurazione del Wi-Fi si ipotizza un \emph{captive portal}; la rete deve poter essere cambiata dalla web app senza ripetere l'associazione.

\textbf{Decisioni prese:}
\begin{itemize}
\tightlist
\item[$\checkmark$] \textbf{[VI-\vnum-13]} Le modalità di associazione del dispositivo (RF43, RF44) e i relativi aspetti di sicurezza vanno chiariti con la proponente prima di modificare i requisiti.
\end{itemize}

\textbf{Incaricati:} Rampazzo\\
\textbf{Scadenza:} 2026-11-20

\hypertarget{odg-4}{%
\subsection{Requisiti di qualità, prestazionali e di vincolo}}
\textbf{Discussione:}\\
I valori dei requisiti prestazionali della bozza sono difficili da giustificare senza conoscere le tecnologie. La portata del microfono resta per ora di 3 metri; la tecnologia della web app verrà decisa in progettazione.

\textbf{Decisioni prese:}
\begin{itemize}
\tightlist
\item[$\checkmark$] \textbf{[VI-\vnum-14]} I valori dei requisiti prestazionali vanno discussi con la proponente prima di essere fissati; RP6 è da rivedere;
\item[$\checkmark$] \textbf{[VI-\vnum-15]} Rimozione dell'esempio dal requisito RV21.
\end{itemize}

\hypertarget{odg-5}{%
\subsection{Diagrammi, casi d'uso e unione delle due bozze}}
\textbf{Discussione:}\\
Esistono due bozze dell'Analisi dei requisiti su \emph{branch} diversi: si prosegue dalla prima, recuperando dalla seconda i diagrammi UML rifatti da Stangherlin. Durante la riunione il gruppo ha lavorato insieme sui casi d'uso del progetto, i cui diagrammi vanno corretti secondo lo standard UML; il diagramma degli stati è ancora una bozza. Il documento non è ancora consegnabile: vanno approfondite le tecnologie del server (MCP, integrazione di Whisper e Piper).

\textbf{Decisioni prese:}
\begin{itemize}
\tightlist
\item[$\checkmark$] \textbf{[VI-\vnum-16]} L'Analisi dei requisiti prosegue dalla bozza di Beltrame e Toffoli, integrando le parti valide della bozza di Stangherlin e Wu;
\item[$\checkmark$] \textbf{[VI-\vnum-17]} I diagrammi dei casi d'uso vanno verificati rispetto allo standard UML presentato a lezione;
\item[$\checkmark$] \textbf{[VI-\vnum-18]} L'Analisi dei requisiti non viene consegnata alla proponente nella riunione del 20 novembre 2026; il lavoro prosegue la settimana successiva.
\end{itemize}

\textbf{Incaricati:} Stangherlin, Wu\\
\textbf{Scadenza:} 2026-11-27

\hypertarget{odg-6}{%
\subsection{Preparazione della riunione con la proponente del 20 novembre 2026}}
\textbf{Discussione:}\\
Alla riunione si porta una lista delle funzionalità da classificare insieme alla proponente; le questioni aperte già risolte vengono rimosse dal documento. Prima della riunione va chiarito se il \emph{firmware} sarà scritto in Rust.

\textbf{Decisioni prese:}
\begin{itemize}
\tightlist
\item[$\checkmark$] \textbf{[VI-\vnum-19]} Alla riunione con la proponente si porta una lista delle funzionalità da classificare per priorità e per consegna (\gl{RTB}, \gl{PoC}, MVP);
\item[$\checkmark$] \textbf{[VI-\vnum-20]} Rimozione delle questioni aperte già risolte e riformulazione di quella sulle aspettative per le consegne.
\end{itemize}

\textbf{Incaricati:} Toffoli, Beltrame\\
\textbf{Scadenza:} 2026-11-20

\hypertarget{varie-ed-eventuali}{%
\subsection{Varie ed eventuali}}
\textbf{Discussione:}\\
Milanese, in qualità di Amministratore, richiama alcune regole sull'uso di GitHub:
\begin{itemize}
\tightlist
\item prima di aprire una \emph{issue} o un \emph{branch} si controlla che non ne esistano già per la stessa attività;
\item nel titolo delle \emph{issue} non si ripete il nome del documento dopo la sigla;
\item chi svolge l'attività apre la propria \emph{issue} e il relativo \emph{branch};
\item una \emph{issue} si chiude solo quando la \emph{pull request} collegata viene unita;
\item il Verificatore raccoglie le correzioni in un unico commento.
\end{itemize}

La verifica della \emph{pull request} del Glossario v0.2 passa a Wu, poiché Salmaso ne è redattore. Il termine \gl{linter} e gli altri termini dell'Analisi dei requisiti vanno aggiunti al Glossario.

\textbf{Decisioni prese:}
\begin{itemize}
\tightlist
\item[$\checkmark$] \textbf{[VI-\vnum-21]} Prima di aprire una \emph{issue} o un \emph{branch} si controlla che non ne esistano già per la stessa attività;
\item[$\checkmark$] \textbf{[VI-\vnum-22]} Il Verificatore riporta le correzioni di una \emph{pull request} in un unico commento.
\end{itemize}

\textbf{Motivazione:}\\
Evitare lavoro duplicato e ridurre i tempi di verifica.

\newpage

% ---------- Azioni da intraprendere ----------
% Una riga per ogni attività assegnata durante la riunione. Priorità: Alta, Media o Bassa.
\hypertarget{azioni-da-intraprendere}{%
\section{Azioni da intraprendere}}

\begin{center}
\begin{tabularx}{\linewidth}{>{\hsize=2.0\hsize\raggedright\arraybackslash}X >{\hsize=1.0\hsize\raggedright\arraybackslash}X >{\hsize=0.5\hsize\centering\arraybackslash}X >{\hsize=0.5\hsize\centering\arraybackslash}X}
\rowcolor{bfblue}\thc{Descrizione} & \thc{Incaricati} & \thc{Scadenza} & \thc{Priorità} \\
Preparare la lista delle funzionalità da discutere con la proponente & Toffoli & 2026-11-20 & Alta \\ \rl
Chiarire la scelta del linguaggio del \emph{firmware} (Rust) & Beltrame, Toffoli & 2026-11-20 & Alta \\ \rl
Chiedere alla proponente associazione del dispositivo, chiamate tra dispositivi, modello locale e requisiti prestazionali & Rampazzo & 2026-11-20 & Alta \\ \rl
Discutere le osservazioni di verifica sulla \textit{Bill of Materials} & Salmaso, Beltrame, Toffoli & 2026-11-20 & Alta \\ \rl
Caricare sul proprio \emph{branch} i diagrammi UML corretti & Stangherlin & 2026-11-19 & Media \\ \rl
Caricare lo script del Glossario e verificarne la \emph{pull request} & Wu & 2026-11-20 & Media \\ \rl
Applicare all'Analisi dei requisiti le modifiche decise e unire le due bozze & Beltrame, Toffoli, Stangherlin, Wu & 2026-11-27 & Alta \\ \rl
Verificare i diagrammi dei casi d'uso rispetto allo standard UML & Stangherlin, Wu & 2026-11-27 & Media \\ \rl
Approfondire l'integrazione di Whisper e Piper nel server & Beltrame, Toffoli & 2026-11-27 & Media \\ \rl
Aggiungere al Glossario \emph{linter} e i termini dell'Analisi dei requisiti & Milanese & 2026-11-27 & Bassa \\ \rl
\end{tabularx}
\end{center}

% ---------- Prossima riunione ----------
\hypertarget{prossima-riunione}{%
\section{Pianificazione della prossima riunione}}

\begin{itemize}
\tightlist
\item \textbf{Data proposta:} 2026-11-20;
\item \textbf{Ora prevista:} 15:00;
\item \textbf{Luogo / modalità:} in remoto, tramite Google Meet, con la proponente Uranion;
\item \textbf{Ordine del giorno preliminare:} \textit{Bill of Materials} e scheda scelta, priorità delle funzionalità per RTB, PoC e MVP, questioni aperte dell'Analisi dei requisiti.
\end{itemize}

% ---------- Firme ----------
% Le firme si aggiungono dopo la verifica: per ciascun firmatario, subito prima di \rule,
% inserire la sua immagine della firma presa da assets/ (file firma_<nome>), racchiusa in
% un makebox largo 0pt, come nei verbali precedenti. Finché le immagini mancano,
% il sito considera il verbale non firmato.
\vspace{35pt}
\noindent
\begin{minipage}[t]{0.45\textwidth}
\textbf{Firma Segretario:}\\[18pt]
\makebox[0pt][l]{\raisebox{0mm}{\includegraphics[width=40mm]{firma_lorenzo.png}}}\rule{55mm}{0.5pt}\\[2pt]
\small Lorenzo Milanese
\end{minipage}
\hfill
\begin{minipage}[t]{0.45\textwidth}
\begin{flushright}
\textbf{Firma Verificatore:}\\[18pt]
\makebox[0pt][l]{\raisebox{0mm}{\includegraphics[width=40mm]{firma_nicholas.png}}}\rule{55mm}{0.5pt}\\[2pt]
\small Nicholas Salmaso
\end{flushright}
\end{minipage}

\end{document}
